% !TeX root = ../../tesis.tex
% §3.1.2 — definiciones operativas (#36, D36-8). La revisión bibliográfica que
% estaba en 3-1-2…3-1-5 pasó a §2.1; aquí quedan definición + párrafo del OE.

% -----------------------------------------------------------------------
\subsection{Conceptos Teórico-Analíticos}
\label{subsec:conceptos-operativos}
% -----------------------------------------------------------------------

\subsubsection{Vulnerabilidad Urbana}
\label{subsec:vulnerabilidad-urbana}

En esta tesis, la vulnerabilidad de una red de transporte es la pérdida de
eficiencia que sufre cuando falla o se retira uno de sus nodos o aristas, como
la formulan \citet{ouyang2009vulnerability}. Otras definiciones del término, y
la diferencia entre fallas aleatorias y ataques dirigidos, se revisan en
§\ref{subsec:ea-vulnerabilidad}.

La fragilidad estructural revisada en §\ref{subsec:ea-vulnerabilidad} es el
sustento teórico del Objetivo Específico~3 (§\ref{subsec:obj-especificos}):
localizar los nodos cuya remoción comprometería la conectividad global del
sistema. Los indicadores de Centralidad de Intermediación $B(v)$ y Robustez
Geométrica $\Delta E$, definidos en §\ref{subsec:centralidad-intermediacion} y
§\ref{subsec:robustez-geometrica}, son la operacionalización computacional de los
conceptos de falla dirigida y disminución de eficiencia.

\subsubsection{Planeación Estratégica}
\label{subsec:planeacion-estrategica}

La planeación estratégica del transporte consiste aquí en probar un cambio de
la red sobre su modelo antes de construirlo: se calculan los mismos indicadores
en la red actual y en la red modificada y se comparan, lo que
\citet{fortin2016innovative} llaman \termino{comparación de redes}. Los
antecedentes de esta forma de planear están en §\ref{subsec:ea-herramientas-datos}.

La comparación cuantitativa de redes es el núcleo metodológico del Objetivo
Específico 5 (§\ref{subsec:obj-especificos}): calcular los indicadores $DI$, $T$
y $B(v)$ sobre la red actual y sobre la red ampliada con el corredor anillar, y
comparar los resultados, es esta forma de planeación estratégica aplicada al
caso de la \zmvm{}. Con ella se evalúa el Anillo Periférico como escenario
hipotético. Los resultados del Capítulo~\ref{ch:analisis-red-anillo} no
anticipan cómo operaría el corredor; muestran qué cambia en la red al agregarlo,
bajo los supuestos que ahí se declaran.

\subsubsection{Eficiencia Urbana}
\label{subsec:eficiencia-urbana}

La eficiencia urbana de una red de transporte es su capacidad para unir
cualquier par de puntos por el camino de menor costo ---en tiempo, transbordos
o distancia--- y para que la población alcance sus destinos en tiempos
razonables. La literatura la mide en tres planos, topológico, operativo y
espacial, que se revisan en §\ref{subsec:ea-indicadores}.

A cada plano le corresponde uno de los indicadores que pide el Objetivo
Específico~2 (§\ref{subsec:obj-especificos}). El Factor de Desviación $DI$ mide
la eficiencia de la topología; el Tiempo de Viaje Promedio $T$ incluye las
condiciones reales de operación, y el Nivel de Cobertura $C$ mide el plano
espacial de la accesibilidad. Calculados juntos sobre el grafo multimodal, los
tres dan el diagnóstico de la red actual de la \zmvm{}.

\subsubsection{Interconectividad entre Redes}
\label{subsec:interconectividad}

La interconectividad entre redes es el grado en que modos y líneas distintos
funcionan como un solo sistema. En el grafo se representa con aristas de
transferencia entre nodos de modos o líneas diferentes, y su calidad se mide en
cada parada (§\ref{subsec:ea-indicadores}).

\label{subsubsec:grafo-multimodal}
La representación matemática de la interconectividad entre redes exige la
construcción de un \termino{grafo multimodal} en el que coexistan, sobre una
misma estructura topológica, los nodos y aristas de los distintos modos de
transporte, vinculados mediante aristas artificiales de transferencia
debidamente ponderadas. \citet{mishra2015} señalan que, a diferencia de un
enlace en una red vial ---que es un segmento físico que conecta un nodo con
otro---, un enlace en una red de transporte público multimodal forma parte de
una línea que sirve una secuencia de paradas; dado que diferentes líneas pueden
servir una misma parada, pueden existir múltiples aristas entre los mismos nodos
del grafo, lo que exige representaciones formales más complejas que las de una
red monomodal.

En el modelo computacional implementado en este trabajo, la red de transporte
metropolitana se representa como un grafo dirigido y ponderado
$\grafo = \{\nodos,\, \aristas,\, W\}$, donde $\nodos$ comprende el conjunto de
todas las paradas de todos los modos, $\aristas$ incluye tanto las aristas de
trayecto (intra-línea) como las aristas de transferencia (inter-modo), y $W$ es
la función de peso que asigna a cada arista un costo compuesto ---tiempo de
viaje, tiempo de espera, tiempo de caminata--- calculado a partir de los datos
\gtfs{} de la red. Sobre este grafo unificado, los algoritmos de caminos mínimos
permiten evaluar la interconectividad del sistema de manera integral: una
transferencia deficiente se traduce en un peso elevado en la arista de
transbordo correspondiente, lo que el algoritmo reflejará automáticamente como
una ruta subóptima, señalando así los nodos de transferencia que requieren
intervención prioritaria.

La construcción del grafo multimodal $\grafo = \{\nodos,\, \aristas,\, W\}$ que
formaliza esta subsección es el resultado tangible que persigue el Objetivo
Específico 1 (§\ref{subsec:obj-especificos}): sistematizar los datos \gtfs{} de
todos los operadores de la \zmvm{} y unificarlos en una representación
matemática que haga posible el análisis cuantitativo. El \texttt{Apímetro}, descrito en
el \autoref{ch:Capas-codigo}, es la herramienta que ejecuta esta construcción de
forma automatizada y replicable.
